% DRAFT (not applied to report.tex). Replaces the "Throughput, not latency" paragraph and the
% "% TODO: update after the professor's answer on latency expectations." line in the Discussion.
% Sources: logs/pace_scan_11nov.txt, logs/pace_sweep.log, logs/pace_sweep_times.csv,
%          logs/pace_{s60,s300,s1500,unpaced}.txt, logs/overhead_clean.log, logs/overhead_clean_times.csv

\subsection{Latency}\label{sec:latency}
\paragraph{Probe and paced replay.} With \texttt{--latency-log}, the parser
stamps each event with the wall-clock time at which it enters the pipeline, and
the window function writes, for every result, the time from that stamp (of the
window's last event) to the moment the result is emitted. All tasks run in one
JVM, so the stamps are comparable. Read as fast as possible, the input is a
backlog, not a stream, so we added a paced replay
(\texttt{--replay-speed}~$S$): right after parsing, each reader holds an event
until $\textit{start} + V/S$ of wall-clock time has passed, where $V$ is the
event time elapsed so far (the largest event time seen, with forward jumps over
10 minutes counted as 10 minutes). The arrival stamp is renewed when the event
is released, so the deliberate wait is not counted as latency. We replay one
weekday without placeholder rows (11 November: 11.8M price events over
15.6\,h of event time, no gap over 10 minutes) with tolerance 0.1\,h at
parallelism 8. The results are identical at every speed.

\begin{table}[t]
\caption{Paced replay of 11 November (tolerance 0.1\,h, parallelism 8,
368,727 results). \emph{Measured}: wall-clock latency of each result.
\emph{Excess}: measured $-$ (built-in wait), where the built-in wait is the
event time from the window's last event to window end $+$ tolerance (or to
the end of the input), divided by $S$. Times in ms.}\label{tab:paced}
\small
\begin{tabular}{lrrrrr}
\toprule
 & $S{=}60$ & $S{=}300$ & $S{=}1500$ & unpaced \\
\midrule
Run time (s)            & 938.7 & 190.7 & 41.2 & 22.7 \\
Events/s                & 12.6k & 62.1k & 287k & 521k \\
\midrule
Measured median         & 6,397 & 1,381 & 406 & 313 \\
Measured p95            & 8,972 & 1,886 & 514 & 466 \\
Measured p99            & 11,051 & 2,347 & 576 & 614 \\
Measured max            & 14,987 & 3,003 & 794 & 810 \\
\midrule
Excess median           & 145 & 120 & 147 & 313 \\
Excess p95              & 257 & 219 & 211 & 466 \\
Excess p99              & 286 & 239 & 272 & 614 \\
\bottomrule
\end{tabular}
\end{table}

\paragraph{Results.} Table~\ref{tab:paced} shows that the measured latency is
almost entirely the wait built into event-time windows. A window can fire only
when event time has passed its end plus the tolerance, so its last event waits
between 6 and 11 minutes of event time (0.1\,h tolerance plus up to one 5-minute
window), which at $S{=}60$ is 6 to 11\,s of wall-clock time. What the pipeline
adds, the excess, stays at a median of 120--150\,ms and a p99 below 290\,ms up
to $S{=}1500$ (287k events/s); this is consistent with Flink's periodic
watermark emission every 200\,ms. Up to $S{=}1500$ every paced run ends about
3.7\,s after its built-in wait ($V/S$), so the pipeline keeps up; only the
unpaced run, at 521k events/s, raises the excess (median 313\,ms). From the
unpaced run time we expect saturation from about $S \approx 3{,}000$.

\paragraph{End of the day.} In the 686 windows that end in the last 30 minutes
before the day's last event, trading thins out, and the event that moves event
time past a window's end plus tolerance can come minutes later. That wait is
still deliberate pacing: counted from the event that actually fires the window,
the excess in these windows is at most 43\,ms ($S{=}60$), 236\,ms ($S{=}300$)
and 116\,ms ($S{=}1500$). Counted to the window end plus tolerance instead, it
reaches 4.1\,s at $S{=}60$ (about 4 minutes of event time).

\paragraph{Limits of the saturated replay.} Our first latency measurement read
the full week as fast as possible (median 40\,s, p99 106\,s at 0.1\,h); there,
events queue in buffers and windows wait for the watermark to sweep through the
backlog, so those numbers reflect the length of the run, not the latency of the
pipeline. A live deployment would see the paced numbers: the tolerance and
window size set the latency, and the pipeline adds a fraction of a second.

\paragraph{Cost of the probe code.} With the probe off (no
\texttt{--latency-log}, no \texttt{--replay-speed}), the probe build still
carries the extra timestamp field through the pipeline. On the full week
(200\,h, parallelism 8, three alternating runs each, idle machine) the main
build took 48.1, 51.5 and 50.4\,s (mean 50.0\,s) and the probe build 52.7,
52.9 and 52.3\,s (mean 52.7\,s, $+5.3\,\%$). The difference (2.7\,s) is smaller
than the spread of the main runs (3.5\,s), but all three probe runs were slower
than all three main runs. The outputs are identical.
